% !TEX program = platex
% 数学ノート集・共通テンプレート v2（積分論2026年度第1部の版面を参考）
% A5判。章内で定義・定理・例・注意等の番号を共通化。
% コンパイル： platex -interaction=nonstopmode math_note_template_v2.tex (2回)
%              dvipdfmx math_note_template_v2.dvi
\documentclass[a5paper,10pt,onecolumn]{jsarticle}
\usepackage[dvipdfmx,top=22mm,bottom=20mm,inner=19mm,outer=18mm,headheight=13pt,headsep=6mm]{geometry}
\usepackage{amsmath,amssymb,amsthm,mathtools}
\usepackage{enumitem,fancyhdr,titlesec}
\usepackage[dvipdfmx,bookmarks=true,colorlinks=false]{hyperref}
\usepackage{pxjahyper}

% ----- 差し替える文書情報 -----
\newcommand{\NoteTitle}{数理モデルの構築と次元解析}
\newcommand{\NoteSeries}{数学ノート集}

% ----- 文字・組版 -----
% 参考PDFの質素な明朝＋Computer Modernの印象を引き継ぎながら，
% A5判で窮屈にならないよう本文の行間を少し確保する．
\linespread{1.13}
\setlength{\parindent}{1zw}
\setlength{\parskip}{0pt}
\setlength{\abovedisplayskip}{6pt plus 2pt minus 2pt}
\setlength{\belowdisplayskip}{6pt plus 2pt minus 2pt}
\setlength{\abovedisplayshortskip}{3pt plus 2pt minus 1pt}
\setlength{\belowdisplayshortskip}{4pt plus 2pt minus 1pt}
\setlist[enumerate]{topsep=3pt,itemsep=1pt,parsep=0pt,leftmargin=2.4zw}
\titleformat{\section}{\normalfont\normalsize\bfseries}{\thesection}{0.8zw}{}
\titlespacing*{\section}{0pt}{2.0ex plus .4ex}{.9ex}
\titleformat{\subsection}{\normalfont\normalsize\bfseries}{\thesubsection}{0.7zw}{}
\titlespacing*{\subsection}{0pt}{1.5ex plus .3ex}{.6ex}
\raggedbottom

% ----- 柱とページ番号：ページ上部、罫線なし -----
\fancypagestyle{mathnote}{%
  \fancyhf{}%
  \fancyhead[L]{\footnotesize \NoteSeries\ （\NoteTitle）}%
  \fancyhead[R]{\footnotesize \thepage}%
  \renewcommand{\headrulewidth}{0pt}%
  \renewcommand{\footrulewidth}{0pt}%
}
\pagestyle{mathnote}

% ----- 節内通し番号の定義・定理群 -----
\newtheoremstyle{mathnotetheorem}
 {4.5pt}{4.5pt}{\normalfont}{}%
 {\bfseries}{}{0.65zw}%
 {\thmname{#1}\thmnumber{ #2}\thmnote{（#3）}}
\theoremstyle{mathnotetheorem}
\newtheorem{theorem}{定理}[section]
\newtheorem{definition}[theorem]{定義}
\newtheorem{lemma}[theorem]{補題}
\newtheorem{proposition}[theorem]{命題}
\newtheorem{corollary}[theorem]{系}
\newtheorem{example}[theorem]{例}
\newtheorem{remark}[theorem]{注意}
\renewcommand{\proofname}{証明}
\renewcommand{\qedsymbol}{$\square$}
\makeatletter
\renewenvironment{proof}[1][\proofname]{%
  \par\pushQED{\qed}\normalfont\topsep4\p@\@plus2\p@\relax
  \trivlist\item[\hskip\labelsep\bfseries #1]\ignorespaces
}{\popQED\endtrivlist\@endpefalse}
\makeatother

% ----- 使うノートで適宜拡張する記号 -----
\newcommand{\R}{\mathbb{R}}
\newcommand{\norm}[1]{\lVert #1\rVert}
\newcommand{\abs}[1]{\lvert #1\rvert}

\begin{document}
\section*{到達目標と前提}
与えられた物理量の次元から次元行列を作り，
その核を求めて無次元量を構成できるようになる．
さらに得られた関係が何を決定し，
何を決定しないのかを具体的な力学モデルで判定する．
前提は線形代数の核，単位系，常微分方程式の初歩である．

\section{数理モデルと仮定の区別}
数学的なモデルは，対象の状態変数，パラメータ，初期条件，
それらの関係を表す方程式によって定められる．
モデルが現実をどこまで説明するかは数学的証明だけからは結論できない．
本稿ではモデルの仮定を明示し，その帰結を定理として扱う．
\begin{definition}
時間区間 $J\subset\mathbb R$，
状態空間 $U\subset\mathbb R^m$，
パラメータ集合 $\Theta\subset\mathbb R^r$ とする．
初期値問題
\[
 \dot x(t)=F(t,x(t),\theta),\qquad x(t_0)=x_0
\]
を状態方程式モデルという．
観測写像 $h:U\times\Theta\to\mathbb R^q$ を与える場合，
$y(t)=h(x(t),\theta)$ を観測量という．
\end{definition}
\begin{proposition}[一意性に必要な正則性]
$F$ が連続で，状態変数について局所Lipschitzであるなら，
与えられた $(t_0,x_0,\theta)$ を通る解は
局所的に存在し，その共通定義域上で一意である．
\end{proposition}
\begin{proof}
パラメータ $\theta$ を固定する．
$F(t,x,\theta)$ は通常の常微分方程式の右辺なので，
Picard--Lindel\"ofの定理を適用すればよい．
解の一意性は，同じ初期点から出る二つの局所解の
一致を積分不等式で示し，延長して得られる．
\end{proof}
\begin{remark}
この命題は解が任意の時間まで存在することを意味しない．
大域存在には不変領域や成長条件等を別に検証する必要がある．
\end{remark}
\section{次元ベクトル}
基本次元を $k$ 種類選び，
物理量 $q_j>0$ の次元を
$d_j=(d_{1j},\ldots,d_{kj})^T\in\mathbb Z^k$ とする．
\begin{definition}
次元行列を
$D=(d_1\ \cdots\ d_n)\in\mathbb Z^{k\times n}$ とする．
基本単位の変更係数
$\lambda=(\lambda_1,\ldots,\lambda_k)\in(\mathbb R_{>0})^k$
の作用を
\[
 (\lambda\cdot q)_j=q_j\prod_{i=1}^k\lambda_i^{d_{ij}}
\]
で定める．
量 $q^v=\prod_{j=1}^nq_j^{v_j}$ が無次元とは
$Dv=0$ となることをいう．
\end{definition}
\begin{lemma}[無次元量の不変性]
$Dv=0$ なら任意の単位変換 $\lambda$ に対して
$(\lambda\cdot q)^v=q^v$ である．
\end{lemma}
\begin{proof}
定義から
\[
 (\lambda\cdot q)^v
 =q^v\prod_{i=1}^k
 \lambda_i^{\sum_jd_{ij}v_j}
 =q^v\prod_i\lambda_i^{(Dv)_i}=q^v
\]
となる．
\end{proof}

\section{Buckinghamの$\Pi$定理}
\begin{lemma}[単位変換の軌道]\label{lem:orbit}
$D$ の階数を $r$ とし，
$v_1,\ldots,v_{n-r}$ を
$\ker(D:\mathbb R^n\to\mathbb R^k)$ の基底とする．
$q,q'\in(\mathbb R_{>0})^n$ に対し，次は同値である．
\begin{enumerate}
\item ある $\lambda\in(\mathbb R_{>0})^k$ が存在して
 $q'=\lambda\cdot q$．
\item 全ての $a=1,\ldots,n-r$ について
 $q^{v_a}=(q')^{v_a}$．
\end{enumerate}
\end{lemma}
\begin{proof}
(1)から(2)は前補題による．
逆に(2)とする．成分ごとに対数を取って
$z=\log q'-\log q\in\mathbb R^n$ と置くと
\[
 v_a^Tz=\log((q')^{v_a})-\log(q^{v_a})=0
\]
である．$v_a$ は $\ker D$ の基底なので
$z\in(\ker D)^\perp=\operatorname{im}D^T$ となる．
従ってある $w\in\mathbb R^k$ が存在して
$z=D^Tw$ と書ける．
$\lambda_i=e^{w_i}$ と置けば
$\log q'_j=\log q_j+\sum_i d_{ij}\log\lambda_i$，
従って $q'=\lambda\cdot q$ である．
\end{proof}
\begin{theorem}[Buckinghamの$\Pi$定理]\label{thm:pi}
$D$ の階数を $r$ とし，
$v_1,\ldots,v_{n-r}$ を $\ker D$ の基底とする．
$\pi_a(q)=q^{v_a}$ と置く．
単位変換に関して不変な任意の写像
$f:(\mathbb R_{>0})^n\to\mathbb R$ は，
ある一意な写像
$\Phi:(\mathbb R_{>0})^{n-r}\to\mathbb R$ を用いて
\[
 f(q)=\Phi(\pi_1(q),\ldots,\pi_{n-r}(q))
\]
と書ける．$f$ が滑らかなら $\Phi$ も滑らかに取れる．
\end{theorem}
\begin{proof}
$\pi=(\pi_1,\ldots,\pi_{n-r})$ とする．
まず $\pi$ は全射である．
行列 $V$ の列を $v_a$ とすると
$V^T:\mathbb R^n\to\mathbb R^{n-r}$ は全射なので，
任意の $\xi\in\mathbb R^{n-r}$ に対し
$V^Tz=\xi$ を解ける．
$q_j=e^{z_j}$ と置けば
$\pi_a(q)=e^{\xi_a}$ を得る．
補題\ref{lem:orbit}から $\pi(q)=\pi(q')$ なら
$q,q'$ は同じ単位変換軌道に属する．
$f$ は軌道上一定であるため，
$\Phi(\pi(q))=f(q)$ と定めても代表元によらない．
$\pi$ の全射性で $\Phi$ は一意である．
$f$ が滑らかな場合は線形写像 $V^T$ の右逆行列
$B:\mathbb R^{n-r}\to\mathbb R^n$ を取り，
\[
 \Phi(s)=f\bigl(\exp(B(\log s))\bigr)
\]
と書ける．ここで対数・指数は成分ごとの写像である．
右辺は滑らかなので結論を得る．
\end{proof}
\begin{example}[単振り子の周期の次元解析]
振り子の周期 $T$ が長さ $\ell>0$ と重力加速度
$g>0$ のみから定まると仮定する．
基本次元を長さ $L$，時間 $T_0$ とすれば
\[
 [T]=T_0,\quad [\ell]=L,\quad [g]=LT_0^{-2}.
\]
無次元量は $T\sqrt{g/\ell}$ だけなので，
次元解析は
\[
 T=C\sqrt{\ell/g}
\]
という形を与えるが，定数 $C$ 自体は決定しない．
\end{example}
\begin{proof}
$T,\ell,g$ の指数を $(a,b,c)$ とすると
無次元性は $b+c=0$，$a-2c=0$ である．
解空間は $(2,-1,1)$ の実数倍なので，
$T^2g/\ell$ が基本無次元量となる．
周期を正とすれば
$T\sqrt{g/\ell}$ も無次元である．
次元の異なる物理的量の関係は単位変換に不変なので，
定理\ref{thm:pi}によってこの値は無次元定数へ帰着する．
$C=2\pi$ などの値は小振幅近似に基づく運動方程式から
別途導く必要がある．
\end{proof}
\begin{remark}
$\Pi$定理は次元的に許容される関数形を制限するが，
物理法則の真偽や経験データへの適合性を証明するものではない．
\end{remark}

\section{パラメータ同定可能性}
\begin{definition}
パラメータ $\theta\in\Theta$ と観測データ
$G(\theta)\in\mathbb R^m$ の対応を考える．
$G(\theta)=G(\theta')$ なら必ず $\theta=\theta'$ となるとき，
パラメータは大域的に識別可能という．
ある $\theta_0$ の近傍でこの条件が成り立つとき，
$\theta_0$ で局所的に識別可能という．
\end{definition}
\begin{theorem}[Jacobi行列による十分条件]\label{thm:ident}
$\Theta\subset\mathbb R^r$ を開集合，
$G:\Theta\to\mathbb R^m$ を $C^1$ 写像とし，
$m\ge r$ とする．ある $\theta_0$ において
$\operatorname{rank}DG(\theta_0)=r$ なら，
$\theta_0$ は局所的に識別可能である．
\end{theorem}
\begin{proof}
階数が $r$ なので $DG(\theta_0)$ の
ある $r$ 行を取り出して得られる $r\times r$ 小行列は可逆である．
対応する出力座標への射影を
$\pi:\mathbb R^m\to\mathbb R^r$ と置くと，
$D(\pi\circ G)(\theta_0)$ は可逆である．
逆写像定理によって $\pi\circ G$ は
$\theta_0$ のある近傍で一対一となる．
その近傍で $G(\theta)=G(\theta')$ なら
$\pi G(\theta)=\pi G(\theta')$ なので
$\theta=\theta'$ である．
\end{proof}
\begin{remark}
Jacobi行列の階数が不足することは識別不能の必要条件ではない．
例えば $\theta\mapsto\theta^3$ は零点で微分が零であるが，
任意の実数上で一対一である．
定理\ref{thm:ident}は十分条件であり必要条件ではない．
\end{remark}
\begin{example}[指数減衰の同定]
正のパラメータ $a,k>0$ に対し
$y(t)=a e^{-kt}$ とする．異なる二時刻
$t_1<t_2$ における正確な観測値
$y(t_1),y(t_2)>0$ から $(a,k)$ は一意に定まる．
\end{example}
\begin{proof}
観測値の比を取ると
$y(t_2)/y(t_1)=e^{-k(t_2-t_1)}$ である．
従って
\[
 k=-{\log(y(t_2)/y(t_1))\over t_2-t_1},
 \qquad a=y(t_1)e^{kt_1}
\]
となる．これらは観測値からただ一つ定まる．
\end{proof}
\section{感染症モデルの一例}
$S,I,R$ をそれぞれ感受性者・感染者・回復者の人口比率とし，
$\beta,\gamma>0$ に対して
\[
 \dot S=-\beta SI,\qquad
 \dot I=\beta SI-\gamma I,\qquad
 \dot R=\gamma I
 \tag{SIR}
\]
を考える．これらの方程式を採用すること自体は
外部の現象に関するモデル化の仮定である．
\begin{theorem}[単体の不変性と大域存在]\label{thm:sir}
$S(0),I(0),R(0)\ge0$ で
$S(0)+I(0)+R(0)=1$ とする．
このとき (SIR) には全ての $t\ge0$ 上でただ一つの解があり，
\[
 S(t),I(t),R(t)\ge0,\qquad S(t)+I(t)+R(t)=1
\]
を満たす．
\end{theorem}
\begin{proof}
右辺は多項式なので局所Lipschitzであり，
初期値問題の局所解は一意に存在する．
三式を加えると
$\frac d{dt}(S+I+R)=0$ なので和は常に1である．
局所解の存在する区間で
\[
 \begin{aligned}
 S(t)&=S(0)\exp\left(-\beta\int_0^t I(s)\,ds\right),\\
 I(t)&=I(0)\exp\left(\int_0^t(\beta S(s)-\gamma)\,ds\right).
 \end{aligned}
\]
が成り立つ．
これは各一階線形微分方程式の積分因子法から得られる．
従って $S,I$ は非負であり，
$R(t)=R(0)+\gamma\int_0^t I(s)\,ds$ も非負である．
和が1なので各成分は $[0,1]$ に含まれ，
解の軌道はコンパクト集合
$\{(S,I,R)\in[0,1]^3:S+I+R=1\}$ にとどまる．
その集合上の右辺の微分は有界なので，
局所存在定理は各点において一様な正の時間幅の延長を与える．
従って有限時刻に極大解が終了することはなく，
全ての $t\ge0$ へ延長できる．
\end{proof}
\begin{proposition}[感染者数の初期増減]
$0<I(0)\le1$ とする．
感染者の初期変化の符号は
$\beta S(0)-\gamma$ の符号に一致する．
\end{proposition}
\begin{proof}
(SIR) の第二式を $t=0$ に適用すると
$\dot I(0)=I(0)(\beta S(0)-\gamma)$ である．
$I(0)>0$ なので符号は同じである．
\end{proof}
\begin{definition}
$S$ を全人口に対する比率とした場合に
$\mathcal R_0=\beta/\gamma$ と定める．
\end{definition}
\begin{corollary}[初期感染増加の閾値]
$I(0)>0$ のとき，
$\dot I(0)>0$ であることと
$\mathcal R_0 S(0)>1$ は同値である．
\end{corollary}
\begin{proof}
前命題の条件
$\beta S(0)-\gamma>0$ を $\gamma>0$ で割り，
$\beta S(0)/\gamma>1$ と書き直せばよい．
\end{proof}

\section{単位の情報から振り子の周期を推定する}
\begin{example}[単振り子の次元解析]
長さ $\ell>0$，重力加速度 $g>0$，
質量 $m>0$，初期振幅 $\theta_0\in(0,\pi)$
をもつ理想的な単振り子の周期 $T$ を考える．
物理的次元を時間 $[\mathsf T]$，長さ $[\mathsf L]$，
質量 $[\mathsf M]$ で表すと，
\[
 [T]=[\mathsf T],\quad [\ell]=[\mathsf L],
 \quad [g]=[\mathsf L][\mathsf T]^{-2},
 \quad[m]=[\mathsf M]
\]
であり，角度 $\theta_0$ は無次元である．
$T,\ell,g,m$ の指数をこの順に並べると，
次元行列は
\[
 D=\begin{pmatrix}
 1&0&-2&0\\0&1&1&0\\0&0&0&1
 \end{pmatrix}
\]
となる（行は順に時間・長さ・質量）．
その核は $(1,-1/2,1/2,0)^T$ で張られる．
従って唯一の独立な次元を持たない組合せは
$T\sqrt{g/\ell}$ であり，次元整合性から
\[
 T=\sqrt{\frac\ell g}\,C(\theta_0)
\]
という形まで制限できる．
ただし次元解析だけで $C$ は決定できない．
\end{example}
\begin{proposition}[小振幅近似で定数が決まる]
小振幅近似 $\sin\theta\simeq\theta$ を用い，
運動を
$\ddot\theta+(g/\ell)\theta=0$
でモデル化するなら，零でない任意の解の基本周期は
\[
 T_{\mathrm{lin}}=2\pi\sqrt{\ell/g}
\]
である．
\end{proposition}
\begin{proof}
$\omega=\sqrt{g/\ell}>0$ と置けば
一般解は
$\theta(t)=A\cos(\omega t)+B\sin(\omega t)$．
$(A,B)\ne(0,0)$ の場合の最小正周期は
$2\pi/\omega$ であり，表示式を得る．
これは次元解析から得た関係に
$C=2\pi$ を与える．
\end{proof}
\begin{remark}
有限振幅の正確な運動方程式は
$\ddot\theta+(g/\ell)\sin\theta=0$ である．
振幅を無視して $C(\theta_0)=2\pi$ と結論することはできない．
次元解析は関数形を制限するが，係数や無次元量への依存性まで
単独で決定するものではない．
\end{remark}
\section*{到達確認問題と解答}
\noindent\textbf{問題．}
小振幅モデルで長さを4倍にしたときの周期を，
重力加速度を変えない場合について求めよ．

\medskip\noindent\textbf{解答．}
$T_{\mathrm{lin}}(\ell)=2\pi\sqrt{\ell/g}$
なので
$T_{\mathrm{lin}}(4\ell)=2T_{\mathrm{lin}}(\ell)$．
質量 $m$ は周期の表示式に現れない．

\end{document}
